% ============================================================================
%  SSEE — Sealed Journal Submission (JCAP-standard, single coherent narrative)
%  Author: Mike Edison Almeida Vallejo
%  Built by the 8-phase SEALED METHOD. Each block verified against
%  VERIFICATION_LEDGER.md before being written here.
%  Bibliography: ssee_unified.bib
% ============================================================================
\documentclass[a4paper,11pt]{article}

\usepackage[utf8]{inputenc}
\usepackage[T1]{fontenc}
\usepackage{amsmath,amssymb,amsthm}
\usepackage{graphicx}
\graphicspath{{../results/figures/}}
\usepackage[margin=2.4cm]{geometry}
\usepackage{booktabs}
\usepackage{xcolor}
\usepackage{hyperref}
\hypersetup{colorlinks=true,linkcolor=blue!55!black,citecolor=teal!60!black,
            urlcolor=blue!55!black}
\usepackage{siunitx}
\usepackage{natbib}
\bibliographystyle{plainnat}

\newcommand{\Msun}{M_\odot}
\newcommand{\hMpc}{h\,\mathrm{Mpc}^{-1}}
\newcommand{\kms}{\mathrm{km\,s^{-1}\,Mpc^{-1}}}
\newcommand{\Mpl}{M_{\mathrm{Pl}}}
\newcommand{\Om}{\Omega_{\mathrm{m}}}

% ----------------------------------------------------------------------------
\title{\bfseries The Golden Ratio Sets the Dark-Energy Equation of State:\\
A Falsifiable, Parameter-Frugal Test Against DESI DR2 and Planck}
\author{Mike Edison Almeida Vallejo\thanks{\texttt{mike.almeida1721@gmail.com}}}
\date{\today}
% ----------------------------------------------------------------------------

\begin{document}
\maketitle

\begin{abstract}
\noindent
We present a minimal-parameter framework for the dark sector in which the
\emph{shape} of late-time cosmic acceleration is fixed algebraically by two
dimensionless constants, the golden ratio $\varphi=(1+\sqrt5)/2$ and $\pi$,
rather than by tuned cosmological parameters. The equation of state follows as
bare ratios of structural invariants, $w_0=-T_r/M_v=-0.840$ and
$w_a=-P_{\rm sc}/K_v=-0.670$, with no fitted amplitude; confronted with the DESI
DR2 $w_0w_a$CDM posterior the agreement is at the $0.05\sigma$ level
($\chi^2_{\rm 2D}=0.080$). A Markov-chain analysis combining DESI BAO with a
Planck-anchored prior yields $H_0=66.55\pm0.44~\kms$. The
effective-field-theory description is canonical---$\alpha_T=\alpha_M=\alpha_B=0$
with $\alpha_K(0)=0.4033$---and the inflationary sector predicts
$n_s=1-\varphi^{-7}=0.9656$ and a tensor-to-scalar ratio $r=\varphi^{-10}=
8.1\times10^{-3}$ that is testable by LiteBIRD and CMB-S4. A Boltzmann
calculation shows the two-$\Omega$ bridge constant MIRA is physically required:
without it the acoustic-peak structure degrades by a factor of $\sim22$ in RMS.
We are explicit about the framework's limits: the absolute vacuum scale enters
as a single measured input, the growth amplitude $S_8$ remains in tension with
weak-lensing surveys, and the dynamical origin of MIRA and the
$\varphi$-dark-matter mass are open problems. The result is a falsifiable,
parameter-frugal description of dark energy whose core predictions rest on
$\varphi$ and $\pi$ alone.
\end{abstract}

% ============================================================================
%  MASTER NOTATION TABLE  (Phase 7 deliverable — only USED constants)
% ============================================================================
\section*{Notation: algebraic constants used in this work}
\label{sec:notation}

All constants below are \emph{dimensionless} and built solely from the golden
ratio $\varphi=(1+\sqrt5)/2$ and $\pi$. Constants defined in earlier SSEE
material but \emph{not used} in this submission (e.g.\ VITA) are deliberately
omitted to keep the register honest.

\begin{center}
\begin{tabular}{llll}
\toprule
Symbol & Definition & Value & Role \\
\midrule
$\varphi$      & $(1+\sqrt5)/2$            & $1.61803$  & Golden ratio (UV seed) \\
$\pi$          & ---                       & $3.14159$  & Circle constant (IR seed) \\
$\Omega$       & $\varphi+\pi$             & $4.75962$  & Stability metric \\
$\beta$ (BIAL) & $(\varphi+\pi)/2$         & $2.37981$  & Base coupling scalar \\
$\mathrm{KAL}$ & $\beta+\pi$               & $5.52140$  & Structural viscosity \\
$\mathrm{AURA}$& $\beta+\varphi$           & $3.99784$  & Screening amplitude \\
$\mathrm{MIRA}$& $\mathrm{AURA}/2$         & $1.99892$  & Two-$\Om$ bridge (value clean; mechanism open) \\
$P_{\mathrm{sc}}$ & $\Omega+\varphi$       & $6.37765$  & Dynamical evolution scalar \\
$K_v$          & $2(\varphi+\pi)$          & $9.51925$  & Structural constraint \\
$T_r$          & $3(\varphi+\beta)$        & $11.99353$ & 3D saturation horizon \\
$M_v$          & $\varphi+\pi+K_v$         & $14.27887$ & Maximal dimensional invariant \\
\bottomrule
\end{tabular}
\end{center}

\noindent Derived background quantities (all dimensionless):
\begin{equation*}
w_0=-\frac{T_r}{M_v}=-0.83996,\qquad
w_a=-\frac{P_{\mathrm{sc}}}{K_v}=-0.67000,\qquad
\Om^{\mathrm{dyn}}=1+w_0=0.16005 .
\end{equation*}

% ============================================================================
%  SECTION SKELETON  (Phase 1 — APPROVED structure)
% ============================================================================

\section{Introduction}
\label{sec:intro}

The standard cosmological model, $\Lambda$CDM, fits the available data with six
free parameters~\citep{Planck2018}, but pays for this success with a dark sector whose composition
and energy scale remain unexplained. The recent DESI DR2 measurements, which
favour a dynamical, evolving dark-energy equation of state over a pure
cosmological constant~\citep{AbdulKarim2025}, sharpen the question: if dark
energy is dynamical, what fixes its evolution? A theory that could
\emph{predict} the equation of state, rather than fit it, would mark genuine
progress.

This paper develops one such proposal. The central hypothesis of the
Structural Self-Energy Expansion (SSEE) framework is that the \emph{shape} of
the dark sector---the equation-of-state ratios, the effective-field-theory
function profile, and the relation between the dynamical and cosmological matter
fractions---is fixed algebraically by two dimensionless constants: the golden
ratio $\varphi=(1+\sqrt5)/2$ and $\pi$. From these, a small set of structural
invariants (Table~\ref{sec:notation}) is built, and the late-time observables
follow as bare ratios with no tuned amplitude.

We are deliberate about what this does and does not claim. SSEE is a
\emph{minimal-parameter} framework, not a zero-parameter one: the absolute scale
of the vacuum enters as a single measured input (Postulate~D,
\S\ref{sec:open}), reducing the count of arbitrary inputs from $\Lambda$CDM's
six to of order three. The dark-sector \emph{shape}, however, carries no free
parameter, and this is what makes the framework falsifiable. Once $\varphi$ and
$\pi$ are fixed, the equation of state, the spectral index, the
tensor-to-scalar ratio, and the EFT slip functions are all determined, and any
of them can refute the model.

The paper is organised as a single coherent argument rather than a collection of
results. Section~\ref{sec:core} constructs the equation of state from the
algebraic invariants and confronts it with DESI DR2, derives the two-$\Omega$
structure that distinguishes the dynamical matter fraction
$\Om^{\mathrm{dyn}}=0.160$ from the cosmological $\Om^{\mathrm{cosm}}=0.320$, and
reports the Hubble-constant inference. Section~\ref{sec:eft} gives the canonical
effective-field-theory description, the inflationary predictions, and the
confrontation with Planck PR4, including a Boltzmann test demonstrating that the
bridge constant MIRA is physically necessary. Section~\ref{sec:growth} treats
the growth of structure honestly, including the residual $S_8$ tension.
Section~\ref{sec:falsifiability} sets out the pre-committed prediction register,
and Section~\ref{sec:open} states the model's measured inputs, exact parameter
count, and principal open problems. We conclude in
Section~\ref{sec:conclusions}.

\section{The algebraic core: from $\varphi,\pi$ to the equation of state}
\label{sec:core}
% PHASE 3 (Bloque 3.1) — SEALED below.

\subsection{Construction of the equation-of-state parameters}
\label{sec:w0wa-construction}
The defining ansatz of SSEE is that the late-time dark-energy equation of
state is fixed---without tuning---by ratios of the dimensionless structural
invariants of Table~\ref{sec:notation}. Concretely,
\begin{equation}
  w_0 \;=\; -\,\frac{T_r}{M_v}
        \;=\; -\,\frac{3(\varphi+\beta)}{\varphi+\pi+K_v}
        \;=\; -0.83996,
  \qquad
  w_a \;=\; -\,\frac{P_{\mathrm{sc}}}{K_v}
        \;=\; -\,\frac{\Omega+\varphi}{2(\varphi+\pi)}
        \;=\; -0.67000 .
  \label{eq:w0wa}
\end{equation}
Both quantities are \emph{bare dimensionless numbers}: they are constructed
entirely from $\varphi$ and $\pi$ and carry no free normalisation, no fitted
amplitude, and no scale input. Under the Chevallier--Polarski--Linder
parametrisation $w(a)=w_0+w_a(1-a)$, Eq.~\eqref{eq:w0wa} corresponds to an
$\alpha$-attractor scalar sector with curvature $\alpha=\varphi^4/3$ (derived
in \S\ref{sec:eft}), so the two CPL numbers are not independent fits but two
projections of a single rolling field.

\subsection{Agreement with DESI DR2}
\label{sec:desi-agreement}
We confront Eq.~\eqref{eq:w0wa} with the DESI DR2 $w_0w_a$CDM posterior
(BAO\,$+$\,CMB\,$+$\,DESY5; \citealp{AbdulKarim2025}),
$w_0^{\rm bf}=-0.827\pm0.060$, $w_a^{\rm bf}=-0.750\pm0.290$, with correlation
$\rho=-0.60$. Because the two parameters are strongly correlated, the
meaningful figure of merit is the \emph{joint} (Mahalanobis) distance rather
than two separate one-dimensional offsets. Writing
$\mathbf{d}=(w_0-w_0^{\rm bf},\,w_a-w_a^{\rm bf})$ and the $2\times2$
covariance $\mathsf{C}$,
\begin{equation}
  \chi^2_{\rm 2D}=\mathbf{d}^{\!\top}\mathsf{C}^{-1}\mathbf{d}=0.080,
  \qquad
  p=0.961,
  \qquad
  \text{equivalent } 0.05\sigma .
  \label{eq:chi2}
\end{equation}
The component offsets are $\Delta w_0=-0.22\sigma$ and $\Delta w_a=+0.28\sigma$;
their joint distance is what reduces to $0.05\sigma$. We stress that this
remarkable agreement is obtained with \emph{zero} hand-tuned dark-energy
parameters: the only inputs are $\varphi$ and $\pi$.

\paragraph{Honest caveat: this is agreement with BAO, not with CMB-only.}
The $0.05\sigma$ figure is specific to the DESI DR2 dynamical-dark-energy
posterior, which is dominated by low-redshift BAO ($z\simeq0.3$--$2.3$), where
dark energy actually drives the expansion. Against a \emph{Planck CMB-only}
$w_0w_a$CDM contour ($w_0=-1.03\pm0.03$, $w_a=0.00$, extrapolated from
$z\simeq1100$) the same $(w_0,w_a)$ sits at $\chi^2_{\rm 2D}=40.2$, i.e.\
$6.0\sigma$; against DESY5 supernovae it is $0.67\sigma$. This spread is the
well-known intrinsic BAO--CMB tension of dynamical-dark-energy models and is
\emph{not} resolved by SSEE---no element of the framework reconciles the
$z\simeq1100$ extrapolation with the $z\lesssim2$ measurement. We report the
DESI agreement as what it is: a clean, parameter-free match to the dataset
that directly probes the dark-energy-dominated epoch.

\subsection{A look-elsewhere test: how special are these two ratios?}
\label{sec:look-elsewhere}
A natural objection to any algebraic coincidence is the look-elsewhere
effect: if one is free to combine $\varphi$ and $\pi$ in arbitrarily many
ways, hitting two target numbers is unremarkable. We address this directly,
and we do so with the \emph{model-faithful} denominator rather than an open
integer search. SSEE does not draw $w_0$ and $w_a$ from an unbounded space of
expressions; it draws them from a \emph{closed dictionary} of
$21$ named structural invariants (the Genesis~5.12 set; Table~\ref{sec:notation}
lists the subset that enters the present analysis), specified as a
fixed set rather than expanded to accommodate the data.\footnote{The dictionary
is the Genesis~5.12 constant set \citep{AlmeidaSSEE2026zenodo}; the
look-elsewhere count below is performed over this fixed set, which is not
enlarged or pruned to improve the match. We make no claim of temporal priority
over the DESI release: the agreement of \S\ref{sec:desi-agreement} is a
parameter-free \emph{postdiction}, and genuinely timestamped predictions are
reserved for unreleased data (\S\ref{sec:falsifiability}). Two pairs of invariants
share a value by construction ($\mathrm{IGNIS}=K_v$,
$\mathrm{SOLAR}=\mathrm{MAR}$); the count below deduplicates by value, so each
distinct number is counted once.} The honest question is therefore: among
\emph{all} ratios $A/B$ of named invariants in this fixed dictionary, how many
land on $|w_0|=0.840$ and on $|w_a|=0.670$?

Enumerating every ordered pair yields $245$ distinct ratios in the interval
$(0,5]$. We then count how many fall within a tolerance of each target. At a
tolerance of $\pm0.001$---the sub-percent precision that actually characterises
the DESI match of \S\ref{sec:desi-agreement}---the result is
\begin{equation}
  |w_0|=0.840:\ \text{1 of 245},\qquad
  |w_a|=0.670:\ \text{1 of 245},
  \label{eq:lookelsewhere}
\end{equation}
each hit being the canonical SSEE assignment itself
($\mathrm{AURA}/\Omega = T_r/M_v$ and $P_{\rm sc}/K_v$, both exact to machine
precision). The density of accidental near-hits is thus $\sim0.4\%$ per target:
within the framework's own vocabulary each equation-of-state parameter is
\emph{essentially unique}. This is the opposite of a dense forest of
coincidences, and it is a stronger statement than a generic
$(a\varphi+b\pi)/(c\varphi+d\pi)$ scan, which by construction inflates the
denominator with expressions the model never uses.

The two ratios are moreover not independent draws. Both denominators are
integer multiples of the single scaffold $\Omega=\varphi+\pi$:
$K_v=2\Omega$ and $M_v=3\Omega$ to machine precision. The two CPL numbers are
two projections of one rolling field (\S\ref{sec:w0wa-construction}), not two
separately fitted constants, so the joint look-elsewhere is governed by a
single shared structure rather than by two free choices. We report this as
evidence of internal economy, not as a claim that the dictionary itself is
derived from first principles---that derivation (the dynamical origin of the
factor-of-two and factor-of-three bridges) remains the open mechanism tracked
in \S\ref{sec:open}.

\subsection{The cosmic matter fraction: the two-$\Om$ structure}
\label{sec:two-omega}
The equation of state~\eqref{eq:w0wa} fixes the \emph{dynamical} matter
fraction through the closure relation of a single rolling field,
\begin{equation}
  \Om^{\mathrm{dyn}} \;=\; 1+w_0 \;=\; 0.16005 ,
  \label{eq:omdyn}
\end{equation}
which is the fraction that controls the late-time acceleration. This value is
notably lower than the matter density inferred from the acoustic peaks. SSEE
relates the two through the structural invariant MIRA,
\begin{equation}
  \Om^{\mathrm{cosm}} \;=\; \mathrm{MIRA}\,\cdot\,\Om^{\mathrm{dyn}}
        \;=\; \frac{\mathrm{AURA}}{2}\,(1+w_0)
        \;=\; 1.99892\times0.16005 \;=\; 0.31993 ,
  \label{eq:omcosm}
\end{equation}
the matter fraction that enters the background expansion $E(z)$, the Poisson
equation, and the CMB. We are deliberately precise about the epistemic status
of MIRA, separating two questions that are too often conflated.

\paragraph{Its value is clean; its mechanism is open.} The \emph{value}
$\mathrm{MIRA}=\mathrm{AURA}/2=1.99892$ is fixed algebraically by $\varphi$ and
$\pi$ with no fit whatsoever---it is not adjusted to reproduce
$\Om^{\mathrm{cosm}}$. The \emph{dynamical mechanism} that would derive this
factor-of-two bridge from the field equations is, by contrast, not yet
established; we track it as an explicit open problem (OP-8,
\S\ref{sec:open}). We therefore present MIRA as a postulated structural
relation with a clean numerical value, \emph{not} as a phenomenological
rescaling tuned to data. Stating this plainly is part of the honest
parameter accounting of \S\ref{sec:open}.

\paragraph{MIRA is physically required, not cosmetic.} Independently of its
ultimate derivation, a Boltzmann analysis shows the relation is not optional.
Running a full CLASS computation with $\Om^{\mathrm{cosm}}=0.320$ places the
first three acoustic peaks at $\ell=220,\,535,\,810$, reproducing the $\Lambda$CDM
template to $1.4\%$ RMS. Removing MIRA and feeding the bare dynamical value
$\Om^{\mathrm{dyn}}=0.160$ into the early-time background shifts all three
peaks rightward by $\sim10\%$ ($\ell=240,\,597,\,922$) and degrades the RMS to
$31.5\%$---a factor of $22$. The sound horizon at recombination demands the
higher matter density; SSEE's two-$\Om$ structure is the statement that the
sector driving late-time acceleration ($\Om^{\mathrm{dyn}}$) and the sector
setting the early-time geometry ($\Om^{\mathrm{cosm}}$) are tied by a single
algebraic constant rather than left as two free numbers.

\section{Bayesian validation: MCMC against DESI DR2}
\label{sec:mcmc}
The algebraic predictions of \S\ref{sec:core} are confronted with data in a
full Markov-chain Monte Carlo analysis. The dark-energy sector contributes no
sampled parameters---$w_0$, $w_a$ and the two-$\Om$ structure are fixed by
Eqs.~\eqref{eq:w0wa}--\eqref{eq:omcosm}---so the chain effectively samples
$H_0$ (and $\Omega_b h^2$) against the DESI DR2 BAO data with a CMB-derived
prior, using the affine-invariant ensemble sampler~\citep{emcee}. We ran $100$
walkers $\times\,25{,}000$ steps (seed $42$).

\paragraph{$H_0$ as one quantity in two stages.} We first anchor the
early-time geometry at the CMB-optimal value
\begin{equation}
  H_0^{\mathrm{anchor}} = 67.068~\kms ,
  \label{eq:h0anchor}
\end{equation}
which is the $H_0$ that minimises the Planck likelihood for the SSEE
background; at this anchor the sound horizon and acoustic scale are both well
within tolerance, $r_d = 146.70~\mathrm{Mpc}$ ($1.51\sigma$) and
$\theta_\ast = 0.59636^\circ$ ($0.69\sigma$). Letting the MCMC then correct
$H_0$ against the low-redshift BAO yields the posterior
\begin{equation}
  H_0^{\mathrm{post}} = 66.553 \pm 0.442~\kms ,
  \qquad
  \Omega_b h^2 = 0.02183 \pm 0.00048 ,
  \label{eq:h0post}
\end{equation}
with $r_d = 147.28~\mathrm{Mpc}$ ($0.72\sigma$) at the posterior value. These
are two stages of the same quantity---a CMB anchor and a BAO-corrected
posterior---not two independent measurements.

\paragraph{The BAO--CMB split is a known feature, reported openly.} The
$\sim1.2\sigma$ gap between Eqs.~\eqref{eq:h0anchor} and~\eqref{eq:h0post} is
the same intrinsic BAO--CMB tension discussed in \S\ref{sec:desi-agreement}:
the late-time BAO prefer a slightly lower $H_0$ than the early-time acoustic
geometry. It is a property of dynamical-dark-energy fits to the current data,
not an artefact of SSEE, and we make it explicit rather than hide it inside a
single quoted number. At the CMB anchor all early-universe observables
($r_d$, $\theta_\ast$) sit below $2\sigma$ simultaneously; the residual split
surfaces only when the BAO-corrected posterior is pushed back into a
CMB-calibrated observable, which is precisely where the literature expects it.
Numerically, evaluating $\theta_\ast$ at $H_0^{\mathrm{post}}$ would return
$0.59417^\circ$ ($5.46\sigma$); we quote this not as a tension but to make the
category error explicit---the late-time expansion rate is not the early-time
acoustic calibration, and conflating the two is precisely what the two-stage
treatment avoids.

\section{The CMB and the effective field theory}
\label{sec:eft}

The equation-of-state construction of \S\ref{sec:core} fixes the
late-time expansion history. To test the framework against the cosmic
microwave background and large-scale structure we must specify its behaviour
at the level of linear perturbations. We do so through an effective field
theory (EFT) of dark energy whose functions are again fixed by $\varphi$ and
$\pi$, with no perturbation-level parameters introduced beyond those already
present in the background.

\subsection{The effective field theory of the dark sector}
\label{sec:eft-functions}
In the EFT-of-dark-energy language \citep{Cheung2008,Gubitosi2013,Bellini2014} the linear
dynamics of a single scalar degree of freedom are controlled by four
functions $\alpha_{\rm K}$ (kineticity), $\alpha_{\rm B}$ (braiding),
$\alpha_{\rm M}$ (Planck-mass run) and $\alpha_{\rm T}$ (tensor speed excess).
SSEE fixes them algebraically:
\begin{equation}
  \label{eq:eft-alphas}
  \alpha_{\rm T}=\alpha_{\rm M}=\alpha_{\rm B}=0,
  \qquad
  \alpha_{\rm K}(0)=3\,\Omega_{\rm DE}\,\Om^{\mathrm{dyn}}=0.40330 ,
\end{equation}
so that gravitational waves propagate at the speed of light
($\alpha_{\rm T}=0$, consistent with GW170817~\citep{GW170817}) and the only non-trivial
function is the kineticity, set by the same dynamical matter fraction
$\Om^{\mathrm{dyn}}=0.16005$ that appears in the background. We cross-checked
$\alpha_{\rm K}(0)$ with the Bellini--Sawicki implementation in \texttt{hi\_class}
\citep{Bellini2015} and recovered $0.4033$ to within $0.005\%$. The associated conformal coupling
of the scalar to matter is
\begin{equation}
  \label{eq:betac}
  \beta_c=-\mathrm{AURA}=-3.9978 ,
\end{equation}
extracted from the requirement that the field's present energy fraction equal
$\Omega_{\rm DE}$; the numerical value $\beta_c\simeq-3.990$ agrees with
$-\mathrm{AURA}$ at the $0.2\%$ level. We flag, as a matter of honesty, that the
identification $\beta_c=-\mathrm{AURA}$ is a near-coincidence rather than a proven
algebraic identity (see \S\ref{sec:open}); the robust statement is the
numerical value $\beta_c\simeq-3.99$, fixed with no free parameters.

\subsection{Inflationary sector: $\alpha$-attractor}
\label{sec:inflation}
The primordial perturbations are described by an $\alpha$-attractor
\citep{Kallosh2013} with curvature
\begin{equation}
  \label{eq:alpha-attractor}
  \alpha=\frac{\varphi^4}{3}=2.2847 .
\end{equation}
For the attractor universality class, $n_s=1-2/N_\ast$ and
$r=12\alpha/N_\ast^2$. Taking the number of $e$-folds to be $N_\ast=2\varphi^7
=58.07$---the unique value in the observationally allowed window $[50,60]$ of
the form $2\varphi^n$---yields the two exact identities
\begin{equation}
  \label{eq:ns-r}
  n_s = 1-\varphi^{-7}=0.96556 ,
  \qquad
  r   = \varphi^{-10}=8.13\times10^{-3} .
\end{equation}
The spectral index sits squarely on the Planck~2018 measurement
($n_s=0.9649\pm0.0042$;~\citealp{Planck2018}); the tensor-to-scalar ratio
$r=\varphi^{-10}$ is a
pre-committed, falsifiable prediction (\S\ref{sec:falsifiability}) well below
the current bound $r<0.036$ and within reach of LiteBIRD and CMB-S4. We note
that, within the framework, $n_s$ and $r$ are postulated of this form and
$\alpha=\varphi^4/3$ then follows as an exact algebraic consequence; we do not
claim an independent derivation of the exponent~$7$ from the potential
(see~\S\ref{sec:open}).

\subsection{Confrontation with Planck PR4}
\label{sec:cmb-confront}
With the background of \S\ref{sec:core} and the EFT functions of
Eq.~\eqref{eq:eft-alphas}, we computed the CMB angular power spectra with the
\texttt{CLASS} Boltzmann code~\citep{Lesgourgues2011} and compared against
Planck PR4~\citep{PlanckPR4} (TT, TE, EE and lensing).
Table~\ref{tab:cmb} reports the reduced $\chi^2$.

\begin{table}[t]
  \centering
  \begin{tabular}{lccc}
    \toprule
    Spectrum & SSEE $\chi^2_r$ & $\Lambda$CDM $\chi^2_r$ & $N$ \\
    \midrule
    TT          & 1.047 & 1.043 & 1971 \\
    TE          & 1.041 & 1.040 & 1967 \\
    EE          & 1.041 & 1.039 & 1967 \\
    $\phi\phi$ (lensing) & 0.837 & 0.757 & 9 \\
    \bottomrule
  \end{tabular}
  \caption{CMB confrontation against Planck PR4. SSEE reproduces the acoustic
  spectra at a quality indistinguishable from $\Lambda$CDM, with acoustic peaks
  at $\ell=220,\,536,\,812$. The combined $\Delta\mathrm{BIC}
  (\text{SSEE}-\Lambda\text{CDM})=-20.8$ over $N=5914$ data points reflects the
  two-parameter economy of the background.}
  \label{tab:cmb}
\end{table}

\paragraph{MIRA is physically required, not cosmetic.} The two-$\Om$ structure
of \S\ref{sec:two-omega} is decisively tested by the CMB. Running the Boltzmann
solver with the cosmological matter fraction $\Om^{\mathrm{cosm}}=0.320$ (i.e.\ with
MIRA) places the first three acoustic peaks at $\ell=220,\,535,\,810$, an RMS
deviation of $1.4\%$ from $\Lambda$CDM. Running it instead with the bare
dynamical fraction $\Om^{\mathrm{dyn}}=0.160$ (i.e.\ \emph{without} MIRA) shifts all
three peaks by ${\sim}10\%$ ($\ell=240,\,597,\,922$) and raises the RMS to
$31.5\%$---a factor of ${\sim}22$ worse. The sound horizon at recombination
therefore \emph{demands} $\Om^{\mathrm{cosm}}\approx0.32$: MIRA is not a fit applied
to the data but a structural feature without which the model fails the CMB
outright. Its numerical value is fixed by $\varphi,\pi$; its dynamical origin
remains an open problem (\S\ref{sec:open}).

\section{Structure growth: an honest tension}
\label{sec:growth}

The same linear EFT predicts the growth of structure, and here we report a
genuine, unresolved tension rather than a success. Solving the growth equation
with the SSEE EFT functions gives a suppressed growth factor relative to
$\Lambda$CDM, $G\equiv D_1^{\rm SSEE}/D_1^{\Lambda\rm CDM}=0.866$, and a
present-day amplitude
\begin{equation}
  \label{eq:s8}
  \sigma_8=0.702\pm0.005 , \qquad S_8=0.725\pm0.005
\end{equation}
in the single-sector treatment. This sits in $\sim2\sigma$--$2.8\sigma$ tension
with weak-lensing measurements~\citep{KiDS2021,DESY32022}. A two-sector treatment, in which
a free-streaming component reduces the effective suppression, raises the
amplitude to $\sigma_8=0.794$, $S_8=0.820$, which remains in $\sim2.6\sigma$
tension. Including baryonic feedback via HMcode-2020~\citep{Mead2021} brings the lensing
amplitude to $S_8=0.758$, formally consistent with DES-Y3 ($0.06\sigma$), but we
treat this as an indicative estimate rather than a clean prediction: the
two-sector mass and the feedback amplitude are calibrated quantities, not
algebraic outputs of $\varphi,\pi$.

\paragraph{The honest statement.} SSEE does not currently resolve the
$S_8$/$\sigma_8$ tension from first principles. The linear prediction
$\sigma_8\simeq0.70$--$0.79$ brackets the measured range but is not pinned down
without calibrated nuisance ingredients, and a full non-linear treatment
requires $N$-body simulations that lie beyond the scope of this work. We report
this as an open challenge (\S\ref{sec:open}) rather than a claimed victory.

\section{Falsifiability and the prediction register}
\label{sec:falsifiability}

The defining property of a parameter-frugal framework is that, once the
algebraic invariants are fixed, the numbers it predicts are not free to move.
A fit can always be improved by adding a parameter; SSEE has none to add. We
therefore commit, in advance, to a register of quantitative predictions, each
attached to the observation that will confirm or refute it. We separate the
register into two tiers: \emph{core} predictions that follow directly from the
dimensionless invariants $\varphi$ and $\pi$, and \emph{extension} predictions
that depend on the phenomenological $\varphi$-dark-matter sector and are
therefore conditional on assumptions still under investigation
(\S\ref{sec:open}).

\subsection{Core predictions (clean from $\varphi,\pi$)}

These quantities carry no fitted coefficient. They are bare ratios of the
golden mean and $\pi$, and they cannot be adjusted to accommodate data.

\begin{table}[h]
\centering
\begin{tabular}{@{}llll@{}}
\toprule
Observable & SSEE value & Algebraic form & Falsifying probe \\
\midrule
$w_0$            & $-0.83996$              & $-T_r/M_v$            & DESI Y3 / Euclid BAO+SNe \\
$w_a$            & $-0.67000$              & $-P_{\rm sc}/K_v$     & DESI Y3 / Euclid BAO+SNe \\
$n_s$            & $0.96556$               & $1-\varphi^{-7}$      & Planck / CMB-S4 \\
$r$              & $8.13\times10^{-3}$     & $\varphi^{-10}$       & LiteBIRD / CMB-S4 \\
$\alpha_K(0)$    & $0.40330$               & $3\,\Omega_{\rm DE}\,\Om^{\mathrm{dyn}}$ & DESI / Euclid growth \\
$\alpha_T,\alpha_M,\alpha_B$ & $0$        & exactly zero          & GW standard sirens, ISW \\
\bottomrule
\end{tabular}
\caption{Core prediction register. Each value is fixed by $\varphi$ and $\pi$
alone, with no free parameter available to absorb a discrepancy. The tensor
ratio $r=\varphi^{-10}$ sits an order of magnitude below the current bound
$r<0.036$ and is the cleanest single falsification target: a detection above
$\sim10^{-2}$, or a robust upper bound below $\sim8\times10^{-3}$, refutes the
inflationary sector.}
\label{tab:predictions-core}
\end{table}

The tensor-to-scalar ratio deserves emphasis. Because $n_s=1-\varphi^{-7}$ and
$r=\varphi^{-10}$ are postulated to be of this form, the consistency relation of
the underlying $\alpha$-attractor closes algebraically: with $N_*=2\varphi^7$
and $\alpha=\varphi^4/3$ one obtains
$r=12\alpha/N_*^2=12(\varphi^4/3)/(2\varphi^7)^2=\varphi^{-10}$ identically. The
prediction is thus internally locked rather than independently adjustable. Its
value, $r\simeq8.1\times10^{-3}$, lies squarely within the reach of LiteBIRD and
CMB-S4, which target $\sigma(r)\sim10^{-3}$. SSEE will be decisively tested by
the next generation of $B$-mode experiments.

The vanishing of the EFT slip functions $\alpha_T=\alpha_M=\alpha_B=0$ is
equally sharp: SSEE predicts a gravitational-wave speed identical to light
($c_T=c$) to all orders in the dark sector, no running of the effective Planck
mass, and no braiding. Any measured departure---an anomalous standard-siren
distance, a time-varying $G_{\rm eff}$, or a non-zero braiding signature in the
ISW--galaxy cross-correlation---falsifies the canonical EFT presented in
\S\ref{sec:eft}.

\section{Honest accounting and outstanding challenges}
\label{sec:open}

A framework is only as credible as its honesty about what it does not yet
explain. We therefore state plainly the inputs the model still requires, the
exact parameter count, and the principal open problems---not as caveats buried
in a footnote, but as the research agenda that the present work both motivates
and constrains.

\subsection{What is measured, not derived: the vacuum scale}

SSEE does not predict the absolute energy density of the dark sector. The
overall scale of the vacuum, $\rho_\Lambda\simeq6\times10^{-10}~\mathrm{J\,
m^{-3}}$, enters as a single measured input, which we label \emph{Postulate~D}.
What the algebra fixes is the \emph{shape} of the dark-energy sector---its
equation-of-state ratios $w_0,w_a$, the two-$\Omega$ structure, and the EFT
function profile---not its amplitude. We make no claim to have solved the
cosmological-constant problem; we have moved its free numbers from six tuned
cosmological parameters down to a single measured normalisation plus the
dimensionless invariants $\varphi$ and $\pi$. This is a reduction in arbitrary
inputs, not their elimination.

\subsection{Parameter count: $\sim3$ versus $6$}

We claim a \emph{minimal-parameter} framework, never a zero-parameter one. The
honest accounting is the following. $\Lambda$CDM carries six fitted parameters
$(\Omega_b h^2,\Omega_c h^2,\theta_*,\tau,n_s,A_s)$. SSEE fixes the dark-sector
shape algebraically from $\varphi$ and $\pi$, leaving as genuine free or measured
inputs essentially the vacuum normalisation (Postulate~D), the baryon
density, and the optical depth $\tau$---of order three effective parameters. The
spectral index $n_s=1-\varphi^{-7}$ and the dark-energy ratios are predicted
rather than fitted, which is what shrinks the count. We state this as
``$\sim3$'' deliberately: the precise reckoning depends on whether one counts the
baryon-density ansatz (OP-1, partially resolved) as derived or measured, and we
prefer to under-promise.

\subsection{Open problems as bridges to future work}

Three open problems are load-bearing and we name them explicitly. They define
the path from the present minimal-parameter framework toward a fully derived
theory.

\begin{itemize}
\item \textbf{The MIRA mechanism (OP-8).} The constant
$\mathrm{MIRA}=(3\varphi+\pi)/4=1.99892$ takes a clean algebraic value and is
\emph{numerically necessary}---the CLASS Boltzmann test of \S\ref{sec:eft} shows
the acoustic peaks degrade by a factor of $\sim22$ in RMS without it. What is
missing is a dynamical mechanism that \emph{derives} this value from an action.
Seven candidate mechanisms have been examined and ruled out; the search
continues. The value is not fitted, but its origin is open.

\item \textbf{The $\varphi$-dark-matter mass (OP-9).} The mass
$m_\varphi\simeq5.60$~eV used in the extension sector is a phenomenological
ansatz. Its arithmetic closes numerically, but it rests on a dimensionally
delicate combination and lacks a first-principles derivation from the field
potential. The growth-sector predictions that would descend from it---notably a
free-streaming step in the matter power spectrum near $k_{\rm fs}\simeq
0.493~\hMpc$, testable by DESI Y3 and Euclid---are deferred to future work and
deliberately excluded from the core register of Table~\ref{tab:predictions-core}
until the ansatz is derived.

\item \textbf{The $\beta_c=-\mathrm{AURA}$ coupling (OP-7).} The conformal
coupling agrees with $-\mathrm{AURA}=-3.998$ to $0.2\%$. We present this as a
near-coincidence to be tested, not as a proven identity, pending a QFT-level
derivation of the role assignments.
\end{itemize}

The growth-of-structure tension of \S\ref{sec:growth} is likewise unresolved:
SSEE does not currently bring $S_8$ into agreement with weak-lensing surveys
from first principles, and the non-linear treatment required to settle it
demands $N$-body simulations beyond this work's scope. We record it as an open
challenge.

None of these gaps touch the algebraic core of
Table~\ref{tab:predictions-core}: the dark-energy equation of state, the
spectral index, the tensor ratio, and the vanishing EFT slip functions stand on
the dimensionless invariants alone and are falsifiable today.

\section{Conclusions}
\label{sec:conclusions}

We have presented a minimal-parameter framework in which the shape of the dark
sector is fixed by the two dimensionless constants $\varphi$ and $\pi$. The
late-time equation of state emerges as bare ratios of structural invariants,
$w_0=-0.840$ and $w_a=-0.670$, and matches the DESI DR2 posterior at
$0.05\sigma$ without a single fitted dark-sector parameter. The
effective-field-theory description is canonical, with vanishing tensor, running,
and braiding functions and $\alpha_K(0)=0.4033$; the inflationary sector
predicts $n_s=0.9656$ and $r=8.1\times10^{-3}$; and a Boltzmann calculation
confirms that the two-$\Omega$ bridge is physically required rather than
cosmetic.

The framework's strength is its frugality and its consequent falsifiability.
Its core predictions---the dark-energy ratios, the spectral index, the tensor
ratio, and the vanishing slip functions---rest on $\varphi$ and $\pi$ alone and
will be tested decisively by DESI Y3, Euclid, LiteBIRD, and CMB-S4 within this
decade. We have been equally explicit about its limits: the vacuum scale is a
measured input, the $S_8$ growth amplitude remains in tension with weak lensing,
and the dynamical origins of the MIRA constant and the $\varphi$-dark-matter
mass are unresolved. These define a concrete research agenda rather than fatal
flaws, and none of them touch the algebraic core.

Whether the deep agreement between a handful of golden-ratio combinations and
the measured dark sector reflects an underlying structural principle or a
remarkable set of coincidences is precisely the question that the prediction
register of \S\ref{sec:falsifiability} is designed to settle. The next
generation of cosmological surveys will provide the answer.

\bibliography{ssee_unified}

\end{document}
